% !TEX program = xelatex
\documentclass[10pt, a4paper]{article}

% ===== Essential Packages =====
\usepackage[T1]{fontenc}          % Better font encoding
\usepackage{lmodern}              % Modern Latin Modern font
\usepackage[UTF8, fontset=none]{ctex}  % Chinese support, no preset fontset
\usepackage{xeCJK}                % XeCJK for better font control

% Set CJK fonts (using fonts available on your system)
\setCJKmainfont{Noto Serif CJK SC}    % Main Chinese font
\setCJKsansfont{Noto Sans CJK SC}     % Sans-serif Chinese font
\setCJKmonofont{Noto Sans Mono CJK SC} % Monospace Chinese font

% ===== Math & Symbols =====
\usepackage{amsmath}              % Advanced math
\usepackage{amssymb}              % Additional math symbols
\usepackage{mathrsfs}             % Ralph Smith's Formal Script
\usepackage{braket}               % Dirac bra-ket notation

% ===== Graphics & Tables =====
\usepackage{graphicx}             % Image support
\usepackage{float}                % Better float control
\usepackage{wrapfig}              % Wrapped figures
\usepackage{tikz}                 % Vector graphics
\usepackage[table]{xcolor}        % Colored tables
\usepackage{array}                % Enhanced tables
\usepackage{multirow}             % Multirow tables
\usepackage{tabularx}             % Adjustable-width tables
\usepackage{booktabs}             % Professional table lines
\usepackage{siunitx}              % SI units in tables
\usepackage{caption}              % Better captions
\usepackage{subcaption}           % Sub-captions
\usepackage{adjustbox}            % Box resizing

% ===== Code Listings =====
\usepackage{listings}             % Code blocks
\usepackage{tcolorbox}            % Colored boxes
\tcbuselibrary{listings, breakable} % Listings in tcolorbox

% ===== Misc =====
\usepackage{lipsum}               % Dummy text
\usepackage{footmisc}             % Better footnotes
\usepackage[margin=1in]{geometry} % Page margins
\usepackage{diagbox}

\author{张世航}
\title{光速测量}

\begin{document}
\maketitle
\section{\(2\pi\)}
\begin{tabular}{|c|c|c|c|c|c|}
		\hline
		\(2\pi\) &第一次&第二次&第三次&A(\(10^{-5}m\))&B(\(10^{-5}m\))\\
        \hline
        \(\Omega(MHz)\)&50.407&50.407&50.407&0&577.35(Hz)\\
        \hline
        近程光程(cm)&12.80&12.81&12.82&5.77&5.77\\
        \hline
        \multirow{9}{*}{远程光程(cm)}&7.28&7.29&7.27&5.77&5.77\\
        \cline{2-6}
        &5.18&5.19&5.20&5.77&5.77\\
        \cline{2-6}
        &61.30&61.29&61.30&3.33&5.77\\
        \cline{2-6}
        &3.72&3.74&3.72&6.67&5.77\\
        \cline{2-6}
        &63.68&63.69&63.68&3.33&5.77\\
        \cline{2-6}
        &3.70&3.72&3.72&6.67&5.77\\
        \cline{2-6}
        &65.52&65.54&65.51&8.82&5.77\\
        \cline{2-6}
        &3.72&3.71&3.71&3.33&5.77\\
        \cline{2-6}
        &72.80&72.83&72.82&8.82&5.77\\
        \hline
        光程差(cm)&274.1 & 274.19 &274.11&/&/\\
        \hline
        光速(\(10^{8}m/s\))&2.76& 2.76& 2.76&/&/\\
        \hline
\end{tabular}\\
我们在表格中计算出了每个分量的不确定度,我们现在计算总的不确定度
不确定度的总的计算公式为
\begin{equation}
    U = \sqrt{(\frac{\partial f}{\partial X_1})^2U_1^2+(\frac{\partial f}{\partial X_2})^2U_2^2+\ldots+(\frac{\partial f}{\partial X_n})^2U_n^2}
\end{equation}
我们先汇总A类不确定度,我们以近程光程为例子计算
\[(\frac{\partial c}{\partial y})^2u_{yA}^2 = (\Omega\times2)^2u_{yA}^2\]
则
\begin{equation}
    u_A = \sqrt{(\frac{\partial c}{\partial y})^2u_{A_y}^2+(\frac{\partial c}{\partial x_1})^2u_{A_1}^2+(\frac{\partial f}{\partial X_2})^2u_{A_2}^2+\ldots+(\frac{\partial c}{\partial x_9})^2u_{A_9}^2}= 19594.71947756519 m/s
\end{equation}
\begin{equation}
    u_B = \sqrt{(\frac{\partial c}{\partial \Omega})^2u_{B_\Omega}^2+(\frac{\partial c}{\partial y})^2u_{B_y}^2+(\frac{\partial c}{\partial x_1})^2u_{B_1}^2+(\frac{\partial f}{\partial X_2})^2u_{B_2}^2+\ldots+(\frac{\partial c}{\partial x_9})^2u_{B_9}^2}= 18406.505411701233 m/s
\end{equation}
\begin{equation}
    u = \sqrt{u_A^2+u_B^2} = 26884.05610906694m/s
\end{equation}

\begin{equation}
    c = (2.76\times10^{8} \pm 26884.05610906694)m/s
\end{equation}

\section{\(\pi\)}

\begin{tabular}{|c|c|c|c|c|c|}
		\hline
		\(\pi\) &第一次&第二次&第三次&A(\(10^{-5}\))m&B(\(10^{-5}\))m\\
        \hline
        \(\Omega(MHz)\)&50.407&50.407&50.407&0&577.35(Hz)\\
        \hline
        近程光程(cm)&11.10&11.10&11.11&3.33&5.77\\
        \hline
        \multirow{9}{*}{远程光程(cm)}&7.71&7.72&7.71&3.33&5.77\\
        \cline{2-6}
        &5.72&5.71&5.73&5.77&5.77\\
        \cline{2-6}
        &29.34&29.35&29.33&5.77&5.77\\
        \cline{2-6}
        &3.58&3.58&3.58&0&5.77\\
        \cline{2-6}
        &34.21&34.22&34.20&5.77&5.77\\
        \cline{2-6}
        &1.72&1.74&1.73&5.77&5.77\\
        \cline{2-6}
        &26.30&26.32&26.31&5.77&5.77\\
        \cline{2-6}
        &3.52&3.52&3.52&0&5.77\\
        \cline{2-6}
        &29.22&29.21&29.23&5.77&5.77\\
        \hline
        光程差(cm)&130.22 &130.27& 130.23&/&/\\
        \hline
        光速(\(10^{8}m/s\))&2.63& 2.63& 2.63&/&/\\
        \hline
\end{tabular}\\
\begin{equation}
    u_A = \sqrt{(\frac{\partial c}{\partial y})^2u_{A_y}^2+(\frac{\partial c}{\partial x_1})^2u_{A_1}^2+(\frac{\partial f}{\partial X_2})^2u_{A_2}^2+\ldots+(\frac{\partial c}{\partial x_9})^2u_{A_9}^2}= 30056.927611154417m/s
\end{equation}
\begin{equation}
    u_B = \sqrt{(\frac{\partial c}{\partial \Omega})^2u_{B_\Omega}^2+(\frac{\partial c}{\partial y})^2u_{B_y}^2+(\frac{\partial c}{\partial x_1})^2u_{B_1}^2+(\frac{\partial f}{\partial X_2})^2u_{B_2}^2+\ldots+(\frac{\partial c}{\partial x_9})^2u_{B_9}^2}= 36812.28076852169 m/s
\end{equation}
\begin{equation}
    u = \sqrt{u_A^2+u_B^2} = 47524.340214280186m/s
\end{equation}

\begin{equation}
    c = (2.63\times10^{8} \pm 47524.340214280186)m/s
\end{equation}

\section{低频仪器}
\begin{tabular}{|c|c|}
\hline
\(2\pi\) & 第一次 \\
\hline
\(\Omega(\text{MHz})\) & 3.426 \\
\hline
近程光程(cm) & 21.22 \\
\hline
\multirow{9}{*}{远程光程(cm)} & 147.60 \\
\hline
& 144.02 \\
\hline
& 143.51\\
\hline
& 143.32 \\
\hline
& 143.50 \\
\hline
& 139.61 \\
\hline
& 4.47 \\
\hline
& 86.51 \\
\hline
& 3.61 \\
\hline
& 94.56 \\
\hline
光程差(cm) & 1029.49 \\
\hline
光速(\(10^{8}\,\text{m/s}\)) & 3.0577912 \\
\hline
\end{tabular}
\end{document}